\documentclass[japanese]{enlp}
\usepackage{enlpbbl}
\usepackage[dvipdfmx]{graphicx}
\usepackage{graphicx}
\usepackage{appendix}
\usepackage{amsmath}
\usepackage{amssymb}
\usepackage{array}
\usepackage{adjustbox}
\usepackage{booktabs}
\usepackage{calc}
\usepackage{caption}
\usepackage{enumitem}
\usepackage{here}
\usepackage{multirow}
\usepackage{subfigure}
\usepackage{url}
\usepackage{varwidth}
\allowdisplaybreaks

\Month{February}
\Year{2021}
\Issuetitle{「\qquad 」特集号}

\jsubject{論文種別（査読有）}

\thesis{学位論文}
\jtitle{ここに論文タイトルここに論文タイトルここに論文タイトルここに論文タイトル}
\author{9999999X\quad 土竜 戦士}
\teacher{
    〇〇 〇〇 准教授 \\
}

\institute{
    愛媛大学 大学院理工学研究科 \\
    理工学専攻 数理情報プログラム
}

\date{2099 年 2 月 吉日 提出}

\jabstract{
ここに概要
}

\begin{document}

\maketitle
\tableofcontents

\section{はじめに}
これはセクションです。

日本語文献の引用\cite{araki-1995}。

英語文献の引用\cite{tackstrom-etal-2013-token}。

\subsection{サブセクション}
これはサブセクションです。

\subsubsection{サブサブセクション}
これはサブサブセクションです。

\section{関連研究}
これはセクションです。

\acknowledgment

謝辞はこちらに記述して下さい。

\bibliographystyle{enlpbbl}
\bibliography{sample}

\begin{publication}
\noindent
\begin{itemize}
\item ここに対外発表
\end{itemize}
\end{publication}

\end{document}

